\SourceLeaf{271}%
applications within a given zone. Some rearrangements that help with this calculation therefore deserve a place here. A rearrangement is in any event \textit{necessary} if we take the standard latitude on the sphere in place of the one on the ellipsoid, that is, if we start from a given value of \(Q\) and calculate the other quantities from it.

If we introduce three auxiliary angles \(\varphi,\zeta,\eta\) by setting
\[
\sin\varphi=e \quad .\ .\ .\ .\ .\ .\ .\ .\ .\ .\ .\ .\ .\ .\ .\quad(13)
\]
\[
\operatorname{tang}\zeta=\operatorname{tang}\varphi\cos P^2 \quad .\ .\ .\ .\ .\ .\ .\ .\ .\ .\quad(14)
\]
\[
\operatorname{tang}\eta=\sin\zeta\operatorname{tang}P \quad .\ .\ .\ .\ .\ .\ .\ .\ .\ .\ .\quad(15)
\]
we obtain
\[
\alpha=\frac{1}{\cos\zeta} \quad .\ .\ .\ .\ .\ .\ .\ .\ .\ .\ .\ .\ .\ .\quad(16)
\]
\[
\sin Q=\cos\zeta\sin P \quad .\ .\ .\ .\ .\ .\ .\ .\ .\ .\ .\ .\quad(17)
\]
\[
\cos\eta\cos Q=\cos P \quad .\ .\ .\ .\ .\ .\ .\ .\ .\ .\ .\ .\ .\quad(18)
\]
\[
\sin\eta=\operatorname{tang}\zeta\operatorname{tang}Q \quad .\ .\ .\ .\ .\ .\ .\ .\ .\ .\quad(19)
\]
\[
\operatorname{tang}\tfrac{1}{2}(P-Q)=\operatorname{tang}\tfrac{1}{2}\zeta\mathord{.}\operatorname{tang}\tfrac{1}{2}\eta \quad .\ .\ .\ .\ .\ .\ .\ .\ .\quad(20)
\]
\[
\sin(2\zeta-\varphi)=e\cos2Q \quad .\ .\ .\ .\ .\ .\ .\ .\ .\ .\ .\ .\quad(21)
\]

Equation 18 follows easily by combining 15 and 17. Then 19 follows by combining 15, 17 and 18, 20 from 17, 18 and 19, and finally 21 from 14 and 17.

The calculation is most precise when we use formulas 14, 15 and 20, with \(P\) given, to determine \(\zeta,\eta,Q\) in that order. When \(Q\) is given, we use equations 21, 19 and 20 to obtain \(\zeta,\eta,P\). One or more of the remaining equations may then be used as a check. If we introduce a fourth auxiliary angle \(\theta\) by the formula
\[
\sin\theta=e\sin P \quad .\ .\ .\ .\ .\ .\ .\ .\ .\ .\ .\ .\ .\ .\quad(22)
\]
we obtain
\[
\cos\varphi=\cos\zeta\cos\eta\cos\theta \quad .\ .\ .\ .\ .\ .\ .\ .\ .\ .\quad(23)
\]
and formulas 9 and 12 take the following forms.
\[
A=\frac{a\cos P}{\alpha\cos\theta\cos Q}=\frac{a\cos\eta}{\alpha\cos\theta}=\frac{a\cos\varphi}{\cos\theta^2}=\frac{a\cos\varphi}{1-ee\sin P^2}
\]
\[
k=\frac{\operatorname{tang}(45^0+\frac{1}{2}P)^\alpha}{\operatorname{tang}(45^0+\frac{1}{2}Q)\operatorname{tang}(45^0+\frac{1}{2}\theta)^{\alpha e}}
\]
